% Part: proof-theory
% Chapter: cut-elimination
% Section: ce-largest

\documentclass[../../../include/open-logic-section]{subfiles}

\begin{document}

\olfileid{pt}{cut}{inv}

\olsection{Removing Largest Cuts}

The proof of \olref[seq][inv]{lem:G3c-invert} shown that if \Log{G3c}
!!{prove}s the conclusion of a logical rule (other than
\RightR{\lexists} and \LeftR{\lforall}), it also !!{prove}s each of
the premises of the rule without increasing the !!{height} of the
proof. We can do better and show that this holds also for $\Log{G3c} +
\CutCS$.

As before, we take the \emph{cut rank} of a \CutCS{} inference to be
the depth of its cut !!{formula}.

\begin{lem}\ollabel{lem:inv-G3c-cut}
If $\Log{G3c} + \CutCS$ !!{prove}s the conclusion of a logical
rule~$R$ other than \RightR{\lexists} or \LeftR{\lforall} using !!a{proof}~$\pi$ then it also !!{prove}s each premise
with !!{proof}~$\pi'$ (or !!{proof}s~$\pi'$, $\pi''$ depending on
whether the rule has one or two premises). Moreover,
(a)~$\pheight{\pi'}\le\pheight{\pi}$ (and
$\pheight{\pi''}\le\pheight{\pi}$) and (b)~the number of \CutCS{}
inferences as well as their ranks is the same in~$\pi'$ (and $\pi''$)
as in~$\pi$.
\end{lem}

\begin{proof}
By inspection of the proofs of
\cref{pt:seq:inv:lem:G3c-invert,pt:seq:inv:lem:invert-quant}. That
proof was by induction on~$\pheight{\pi}$. In the base case, we
replaced an axiom by another axiom. So conditions (a) and~(b) are
trivially satisfied. If the last rule of~$\pi$ was~$R$, its immediate
sub-!!{proof}s are the !!{proof}s~$\pi_i$ we want, and (a) and~(b) are
both satisfied. In all other cases, we applied the inductive
hypothesis to the immediate sub-!!{proof}s of~$\pi$, i.e., those
leading to the premise(s) of the last inference~$R$, and then applied
that same rule $R$ to the results. The conditions (a) and~(b) hold for
the immediate sub-!!{proof}s of the new proof, and so (a) and~(b) hold
for the entire proof. It remains to verify the case where the last
inference is~\CutCS. We do this again for one example only: suppose
$\pi$ is !!a{proof} of the conclusion $!B \land !C, \Gamma \Sequent
\Delta$ of \LeftR{\land}, and ends in~\CutCS:
\[
\AxiomC{}
\RightLabel{$\pi_1$}
\Deduce$!B \land !C, \Gamma \fCenter \Delta, !A$
\AxiomC{}
\RightLabel{$\pi_2$}
\Deduce$!A, !B \land !C, \Gamma \fCenter \Delta$
\RightLabel{\CutCS}
\BinaryInf$!B \land !C, \Gamma \fCenter \Delta$
\DisplayProof
\]
The induction hypothesis applies to $\pi_1$ and~$\pi_2$ (which are
also !!{proof}s of instances of the conclusion of $\LeftR{\land}$) and
yields !!{proof}s $\pi_1'$ and~$\pi_2'$ of $!B, !C, \Gamma \fCenter
\Delta, !A$ and $!A, !B, !C, \Gamma \fCenter \Delta$,
respectively. We combine these using \CutCS{} to obtain~$\pi'$:
\[
\AxiomC{}
\RightLabel{$\pi_1'$}
\Deduce$!B, !C, \Gamma \fCenter \Delta, !A$
\AxiomC{}
\RightLabel{$\pi_2'$}
\Deduce$!A, !B, !C, \Gamma \fCenter \Delta$
\RightLabel{\CutCS}
\BinaryInf$!B, !C, \Gamma \fCenter \Delta$
\DisplayProof
\]
Clearly, (a) and~(b) are satisfied.
\end{proof}

Note that this would not work if we used \Cut{} instead of \CutCS, as
then the conclusion of the last !!{proof} would have two copies of
$!B, !C, \Gamma$ in the antecedent and two copies of $\Delta$ in the
succedent. We would have to appeal to the admissibility of
contraction, but the proof of \olref[seq][inv]{prop:G3c-cont-adm} uses
the invertibility lemma. We'd thus be arguing in a circle.

We can use \cref{pt:cut:inv:lem:inv-G3c-cut} to give an alternative
proof of cut elimination. In this proof we do not successively remove
topmost occurrences of \CutCS{} inferences, but occurrences of \CutCS{}
inferences of maximal rank. That is, we start with !!a{proof}~$\pi$ in
$\Log{G3c} + \CutCS$ and remove a cut anywhere in it (possibly
replacing it with cuts of smaller rank)---and then keep doing that
until no cuts remain. The only condition is that above the cut we deal
with, no cuts of the same or higher rank occur.

\begin{lem}\ollabel{lem:max-cut-red-G3c} If $\pi$ is !!a{proof} in
$\Log{G3c}+\CutCS$ ending in a~\CutCS{} inference of rank~$n$ and
otherwise using only rules of~$\Log{G3c}$ and \CutCS{} inferences of
rank $<n$, then there is a proof~$\pi'$ of the same end-sequent
in~$\Log{G3c} + \CutCS$ where every \CutCS{} inference has rank~$<n$.
\end{lem}

\begin{proof}
The !!{proof}~$\pi$ ends in:
\[
\AxiomC{}
\RightLabel{$\pi_1$}
\Deduce$\Gamma \fCenter \Delta, !A$
\AxiomC{}
\RightLabel{$\pi_2$}
\Deduce$!A, \Gamma \fCenter \Delta$
\RightLabel{\CutCS}
\BinaryInf$\Gamma \fCenter \Delta$
\DisplayProof
\]
We distinguish cases according to the form of~$!A$.

Suppose $!A$ is atomic. Observe that by the condition that all
\CutCS{} inferences in~$\pi_1$ and $\pi_2$ must have rank $<
\depth{!A} = 0$, there can be no \CutCS{} inferences in~$\pi_1$
or~$\pi_2$. We show that $\Gamma \Sequent \Delta$ has a cut-free proof
by induction on $\pheight{\pi_2}$. If $\pheight{\pi_2} = 0$, then $!A,
\Gamma \Sequent \Delta$ is an axiom. If $!A$ is not principal, some
atomic $!C$ is !!a{element} of both~$\Gamma$ and $\Delta$, and $\Gamma
\Sequent \Delta$ itself is an axiom. If $!A$ is principal, then
$\Delta = \Delta, !A$. In this case, $\pi_1$ ends in $\Gamma \Sequent
\Delta', !A, !A$. We obtain a cut-free proof of $\Gamma \Sequent
\Delta', !A$ by applying \olref[seq][inv]{prop:G3c-cont-adm}
(admissibility of contraction). If $\pheight{\pi_2} > 0$ it ends in a
rule~$R$. Take a premise of that rule: it has the form $!A, \Gamma',
\Pi \Sequent \Delta', \Lambda$, where $\Pi$ and $\Lambda$ are the
active !!{formula}s of~$R$. \olref[seq][adm]{prop:weak-G3c-adm}
applied to~$\pi_1$ and $\pi_2$ yields !!{proof}s of $\Gamma, \Gamma',
\Pi \Sequent \Delta, \Delta', \Lambda, !A$ and $!A, \Gamma, \Gamma',
\Pi \Sequent \Delta, \Delta', \Lambda$, to which the induction
hypothesis applies. This gives us !!a{proof} of $\Gamma, \Gamma', \Pi
\Sequent \Delta, \Delta', \Lambda$, for each premise of~$R$. Applying
$R$ yields $\Gamma, \Gamma \Sequent \Delta, \Delta$, and
\olref[seq][inv]{prop:G3c-cont-adm} provides the cut-free proof of
$\Gamma \Sequent \Delta$.

If $!A$ is not atomic, we distinguish cases according to the !!{main
operator} of~$!A$. In each case, we apply \olref{lem:inv-G3c-cut} to
obtain !!{proof}s of the premises of the left- and right-rules with
$!A$ as principal formula. We then proceed as in case~E of the proof
of \olref[top]{lem:cut-adm-G3c}.
\begin{enumerate}
  \item $!A \ident \lnot !B$. The !!{proof}s $\pi_1$ and $\pi_2$ end
  in $\Gamma \Sequent \Delta, \lnot !B$ and $\lnot !B, \Gamma \Sequent
  \Delta$. By \olref{lem:inv-G3c-cut}, there are !!{proof}s $\pi_1'$
  and $\pi_2'$ of $!B, \Gamma \Sequent
  \Delta$ and $\Gamma \Sequent \Delta, !B$. The !!{proof}~$\pi'$ is:
  \[
  \AxiomC{}
  \RightLabel{$\pi_2'$}
  \Deduce$\Gamma \fCenter \Delta, !B$
  \AxiomC{}
  \RightLabel{$\pi_1'$}
  \Deduce$!B, \Gamma \fCenter \Delta$
  \RightLabel{\CutCS}
  \BinaryInf$\Gamma \fCenter \Delta$
  \DisplayProof
  \]
  \item $!A \ident (!B \lor !C)$. The !!{proof}s $\pi_1$ and $\pi_2$
  end in $\Gamma \Sequent \Delta, !B \lor !C$ and $!B \lor !C, \Gamma
  \Sequent \Delta$. By \olref{lem:inv-G3c-cut} (inversion for for
  \RightR{\lor} applied to~$\pi_1$), there is !!a{proof} $\pi_1'$ of
  $\Gamma \Sequent \Delta, !B, !C$. By \olref{lem:inv-G3c-cut}
  (inversion for \LeftR{\lor} applied to~$\pi_2$), there is !!a{proof}
  $\pi_2'$ of $!B, \Gamma \Sequent \Delta$ and !!a{proof} $\pi_2''$ of
  $!C, \Gamma \Sequent \Delta$. By applying
  \olref[seq][adm]{prop:weak-G3c-adm} to $\pi_2''$ we also get
  !!a{proof}~$\pi_2'''$ of $!C, \Gamma \Sequent \Delta, !B$. The
  !!{proof}~$\pi'$ is:
  \[
  \AxiomC{}
  \RightLabel{$\pi_1'$}
  \Deduce$\Gamma \fCenter \Delta, !B, !C$
  \AxiomC{}
  \RightLabel{$\pi_2'''$}
  \Deduce$!C, \Gamma \fCenter \Delta, !B$
  \RightLabel{\CutCS}
  \BinaryInf$\Gamma \fCenter \Delta, !B$
  \AxiomC{}
  \RightLabel{$\pi_2'$}
  \Deduce$!B, \Gamma \fCenter \Delta$
  \RightLabel{\CutCS}
  \BinaryInf$\Gamma \fCenter \Delta$
  \DisplayProof
  \]
  \item $!A \ident (!B \land !C)$. Exercise.
  \item $!A \ident (!B \lif !C)$. The !!{proof}s $\pi_1$ and $\pi_2$
  end in $\Gamma \Sequent \Delta, !B \lif !C$ and $!B \lif !C, \Gamma
  \Sequent \Delta$. By \olref{lem:inv-G3c-cut} (inversion for
  \RightR{\lif} applied to~$\pi_1$), there is !!a{proof} $\pi_1'$ of
  $!B, \Gamma \Sequent \Delta, !C$. By \olref{lem:inv-G3c-cut}
  (inversion for \LeftR{\lif} applied to~$\pi_2$), there is
  !!a{proof} $\pi_2'$ of $\Gamma \Sequent \Delta, !B$ and !!a{proof}
  $\pi_2''$ of $!C, \Gamma \Sequent \Delta$. By applying
  \olref[seq][adm]{prop:weak-G3c-adm} to $\pi_2''$ we also get
  !!a{proof}~$\pi_2'''$ of $!C, !B, \Gamma \Sequent \Delta$. The
  !!{proof}~$\pi'$ is:
  \[
  \AxiomC{}
  \RightLabel{$\pi_2'$}
  \Deduce$\Gamma \fCenter \Delta, !B$
  \AxiomC{}
  \RightLabel{$\pi_1'$}
  \Deduce$!B, \Gamma \fCenter \Delta, !C$
  \AxiomC{}
  \RightLabel{$\pi_2'''$}
  \Deduce$!C, !B, \Gamma \fCenter \Delta$
  \RightLabel{\CutCS}
  \BinaryInf$!B, \Gamma \fCenter \Delta$
  \RightLabel{\CutCS}
  \BinaryInf$\Gamma \fCenter \Delta$
  \DisplayProof
  \]
  \item $!A \ident \lforall[x][!B(x)]$. The !!{proof}s $\pi_1$ and
  $\pi_2$ end in $\Gamma \Sequent \Delta, \lforall[x][!B(x)]$ and
  $\lforall[x][!B(x)], \Gamma \Sequent \Delta$. By
  \olref{lem:inv-G3c-cut}, there is !!a{proof}~$\pi_1'$ of $\Gamma
  \Sequent \Delta, !B(c)$.

  Now consider the !!{proof}~$\pi_2$. It ends in $\lforall[x][!B(x)],
  \Gamma \Sequent \Delta$. Moving upward from the end-sequent
  of~$\pi_2$, mark every occurrence of $\lforall[x][!B(x)]$ on the
  left of a sequent that ``leads to'' the occurrence of
  $\lforall[x][!B(x)]$ in the end-sequent of~$\pi_2$, that is:
  (a)~Mark $\lforall[x][!B(x)]$ in the end-sequent of~$\pi_2$. (b)~If
  $\lforall[x][!B(x)]$ occurs in the context of the conclusion of a
  rule and is marked, mark a corresponding occurrence in the
  premise(s) of the rule. (c)~If $\lforall[x][!B(x)]$ is principal in
  the conclusion of a \LeftR{\lforall} rule and is marked, mark the
  corresponding active occurrences of $\lforall[x][!B(x)]$ and $!B(t)$
  in the premise.
  
  Now consider all these marked occurrences of $\lforall[x][!B(x)]$.
  None of them are active in a rule other than \LeftR{\lforall} (only
  active !!{formula}s of \LeftR{\lforall} get marked). Delete all
  marked occurrences of $\lforall[x][!B(x)]$ in $\pi_2$.
  If this is a correct !!{proof}, the marked occurrences of
  $\lforall[x][!B(x)]$ were never active; they all came directly from
  axioms. The !!{proof} ends in~$\Gamma \Sequent \Delta$ and we can
  replace $\pi$ by it. We have a removed a cut of maximal rank.

  Otherwise, what we have is not a correct proof because we have
  deleted !!{formula}s $\lforall[x][!B(x)]$ which were active in some
  \LeftR{\lforall} inferences. We have turned these into
  ``inferences'' of the form
  \[
  \AxiomC{}
  \RightLabel{$\pi_t$}
  \Deduce$!B(t), \Pi \fCenter \Lambda$
  \RightLabel{\LeftR{\lforall}}
  \UnaryInf$\Pi \fCenter \Lambda$
  \DisplayProof
  \]
  Replace each topmost such ``inference'' by
  \[
  \AxiomC{}
  \RightLabel{$\Subst{\pi_1'}{t}{c}[\Pi \Sequent \Lambda]$}
  \Deduce$\Gamma, \Pi \fCenter \Delta, \Lambda, !B(t)$
  \AxiomC{}
  \RightLabel{$\pi_t[\Gamma \Sequent \Delta]$}
  \Deduce$!B(t), \Gamma, \Pi \fCenter \Delta, \Lambda$
  \RightLabel{\CutCS}
  \BinaryInf$\Gamma, \Pi \fCenter \Delta, \Lambda$
  \DisplayProof
  \]
  and add $\Gamma$ to the left and $\Delta$ to the right of every
  sequent below it. Repeat until all \LeftR{\lforall} ``inferences''
  with a marked $!B(t)$ in the premise have been replaced by a
  \CutCS{} inference on some~$!B(t)$. The end-sequent of the resulting
  !!{proof} (and it now is a correct !!{proof}) is of the form
  $\Gamma, \dots, \Gamma \Sequent \Delta, \dots, \Delta$. Apply
  admissibility of contraction to turn it into !!a{proof} of~$\Gamma
  \Sequent \Delta$ and replace $\pi$ with it. We have replaced one cut
  on $\lforall[x][!B(x)]$ with (possibly many) cuts on $!B(t)$, but
  these are all of lower rank. Any cuts already occurring in $\pi_1$ or
  $\pi_2$ remain present, but they are also all of lower rank.
\end{enumerate}
\end{proof}

\begin{prob}
Verify the case where $!A \ident (!B \land !C)$ in the proof of
\olref{lem:inv-G3c-cut}. That is, show that if there is
!!a{proof}~$\pi$ ending in 
\[
\AxiomC{}
\RightLabel{$\pi_1$}
\Deduce$\Gamma \fCenter \Delta, !B \land !C$
\AxiomC{}
\RightLabel{$\pi_2$}
\Deduce$!B \land !C, \Gamma \fCenter \Delta$
\RightLabel{\CutCS}
\BinaryInf$\Gamma \fCenter \Delta$
\DisplayProof
\]
with $\pi_1$ and $\pi_2$ containing only cuts of rank $<\depth{!B
\land !C}$ then there is !!a{proof}~$\pi'$ of $\Gamma \Sequent \Delta$
containing only cuts of rank $<\depth{!B \land !C}$.
\end{prob}

% \begin{prob}
% Verify the case where $!A \ident \lexists[x][!B(x)]$ in the proof of
% \olref{lem:inv-G3c-cut}.
% \end{prob}


\olref[top]{lem:cut-adm-G3c} establishes that we can replace a \CutCS{}
inference of highest rank in !!a{proof} by \CutCS{} inferences of
lower rank. We can again use this lemma to show that we can eliminate
any number of \CutCS{} inferences from !!a{proof} by successively
applying it to the topmost sub-!!{proof}s ending in maximal~\CutCS.

\begin{thm}\ollabel{thm:cut-elim-G3c-largest}
Any !!{proof} $\pi$ in $\Log{G3c} + \CutCS$ can be transformed into a proof in~\Log{G3c}.
\end{thm}

\begin{proof}
  First we prove that all cuts of maximal rank in~$\pi$ can be
  removed. Let $d$ be the maximal rank of cuts occuring in~$\pi$, and
  let $n$ be the number of cuts of rank~$d$. The proof is by induction
  on~$n$. If $n=0$ there is nothing to prove. If $n>0$, pick a topmost
  cut of rank~$d$, and let $\pi_1$ be the sub-!!{proof} ending in it.
  By \olref{lem:max-cut-red-G3c}, there is !!a{proof}~$\pi_1'$ of the
  same end-sequent with all cuts of rank $<d$. Replace $\pi_1$ by
  $\pi_1'$ in~$\pi$. This results in !!a{proof} containing $n-1$ cuts
  of rank~$d$, and the induction hypothesis applies.

  Now the result follows by induction on~$d$. If $d=0$, all cuts are
  atomic, and by the previous result can all be removed. If $d>0$,
  the previous result yields a proof in which all cuts of rank $d$
  have been replaced by cuts of rank~$<d$. Since $d$ was the maximum
  rank of cuts in~$\pi$, we have proof in which the maximal rank of
  cuts is $<d$, and the induction hypothesis applies.
\end{proof}

\end{document}